\subsection{讨论}\label{sec:2.9}
\paragraph{\autoref{sec:2.1}}
关于 Wasserstein 大数律的更详细讨论，见\cite[第11章]{ref123}. 缓慢的收敛速率体现了如下事实：$W_1$ 收敛自动蕴含所有 Lipschitz 检验函数的收敛；并且对 $p>1$，还蕴含所有高阶矩的收敛，如\autoref{prop:1.5}所示. 因此，获得如此强的控制需要大量样本，并不令人意外. Wasserstein 大数律对最优量化及 $k$ 均值聚类的含义，最早在文献\cite{ref259}中于统计背景下展开.

\paragraph{\autoref{sec:2.2}}
Boissard 与 Le Gouic\cite{ref45}最早强调了二进划分论证在控制 Wasserstein 距离中的用途；关于向无界情形的推广，另见\cite{ref108,ref136}. 文献\cite{ref239}进一步讨论了该方法的历史，本节采用的论证版本即取自其中. 容易扩展这一论证，说明收敛速率依赖测度 $\mu$ 的内在维数，而不是环境空间的维数.

二进划分论证也适用于 $p>1$ 的情形，并给出
\begin{equation}\label{eq:2.21}
\mathbb EW_p(\mu_n,\mu)\lesssim_p\sqrt d\cdot\begin{cases}
n^{-1/(2p)},&d<2p,\\
(\log n)^{1/p}/n^{1/(2p)},&d=2p,\\
n^{-1/d},&d>2p.
\end{cases}
\end{equation}
除 $d=2p$ 时的对数因子外，这个速率基本上是尖锐的. 另一方面，本章给出的对偶界不易推广到 $p>1$，因为当 $p>1$ 时，$W_p(\mu_n,\mu)$ 的对偶表述不能产生经验过程.

由三角不等式，$W_p(\mu_n,\nu)$ 向 $W_p(\mu,\nu)$ 的收敛速率，可从 $W_p(\mu_n,\mu)$ 的相应速率推出；但这些速率可能并不尖锐. 举例来说，文献\cite{ref85}证明
\begin{equation}\label{eq:2.22}
\mathbb E|W_2^2(\mu,\nu_n)-W_2^2(\mu,\nu)|\lesssim\begin{cases}
n^{-1/2},&d<4,\\
(\log n)/n^{1/2},&d=4,\\
n^{-2/d},&d>4.
\end{cases}
\end{equation}
注意，当 $\mu\ne\nu$ 时，这个界强于从 \eqref{eq:2.21} 能推出的结果. 其他 $W_p$ 距离也有类似现象\cite{ref220}.

更显著的是，可以证明 $W_p(\mu_n,\nu)$ 向 $W_p(\mu,\nu)$ 的收敛速率，取决于 $\mu$ 与 $\nu$ 的内在维数中较小的一个，见\cite{ref165}. 特别地，若 $\nu$ 支撑在有限个点上，即所谓半离散\index{semi-discrete optimal transport@半离散最优传输（semi-discrete optimal transport）|hyperpage}最优传输问题（semi-discrete optimal transport problem），则 $W_p(\mu_n,\nu)$ 以不受维数灾难影响的速率收敛于 $W_p(\mu,\nu)$. 仅凭 $W_p(\mu_n,\mu)$ 的界无法推出这一事实.

\paragraph{\autoref{sec:2.3}}
链式法的思想隐含地可追溯至 Kolmogorov. \autoref{prop:2.6}的形式称为 Dudley 熵积分（Dudley's entropy integral）\cite{ref120}. 关于链式法及其应用的众多参考资料，可参见\cite{ref122,ref299,ref309,ref310,ref314,ref318}.

\paragraph{\autoref{sec:2.4}}
如前所述，本节结果归于文献\cite{ref5}，这里的论证取自\cite{ref42}. 用 Fourier 变换控制匹配成本的想法最初见于\cite{ref90}，Talagrand\cite{ref299}对其作了广泛发展. 更一般地，所谓匹配问题有大量文献，例如\cite{ref200,ref202,ref299}；近来还利用偏微分方程技术，在 $d=2$ 时导出了非常精细的结果，参见\cite{ref19}.

\paragraph{\autoref{sec:2.5}}
Wasserstein 距离在检验中的应用可见\cite{ref105,ref152,ref161,ref237}. 文献\cite{ref27}研究了以 Wasserstein 距离定义的拟合优度和双样本检验问题；对每个问题，它们证明最优分离量的阶为 $n^{-c/d}$，其中某个常数 $c>1$.

文献\cite{ref241}证明，对于上述拟合优度问题的一个“稳健”版本，在忽略对数因子后，分离量 $n^{-1/d}$ 是必要的；相应假设为
$$H^0:W_1(\mu,\mu^0)<\varepsilon\qquad\text{对}\qquad H_1:W_1(\mu,\mu^0)>2\varepsilon.$$
这个下界可用来得到估计 Wasserstein 距离时几乎尖锐的极小极大下界.

\paragraph{\autoref{sec:2.6}}
\autoref{thm:2.14}的下界来自\cite{ref121}. \autoref{thm:2.15}的极小极大下界由 Singh 与 Póczos\cite{ref286}首次证明. 关于极小极大下界技术的介绍，参见\cite{ref269,ref302,ref318}.

\paragraph{\autoref{sec:2.7}}
Wasserstein 距离下光滑密度的极小极大估计，针对 $W_1$ 的研究见\cite{ref208,ref287}，针对 $W_p$（$p>1$）的研究见\cite{ref240}. $p>1$ 的情形表现不同于 $p=1$：文献\cite{ref240}表明，只有额外假设 $\mu$ 的密度有正的下界，才能在 $p>1$ 时达到\autoref{thm:2.18}的速率；没有这个假设，估计速率严格更差. 关于 $p>1$ 的结果仅限于属于 Besov 类的密度；将文献\cite{ref240}的论证推广到其他密度类，是一个开放问题. 文献\cite{ref111}研究了流形上这个问题的一个版本.

非参数密度估计本身是统计学中的经典课题，只是通常以其他距离度量来研究\cite{ref302}.

\paragraph{\autoref{sec:2.8}}
值得提及的是，IPM 在生成对抗网络（generative adversarial network，GAN），特别是 Wasserstein GAN\cite{ref22}中受到广泛关注；在其中，$\mathcal F$ 被选为一族深度神经网络. 关于 IPM 与 GAN 的统计分析，参见\cite{ref208,ref304}.

最大均值差异最早由文献\cite{ref50}发展，如今已有大量研究，见\cite{ref232}. \autoref{thm:2.22}的速率是广为人知的结果. MMD 的一个重要特例是能量距离（energy distance）类，关于其介绍和参考文献，推荐\cite[第1.2.4小节]{ref142}.

平滑 Wasserstein 距离的良好统计性质最早记录于\cite{ref150}. \autoref{thm:2.24}的简单证明取自\cite{ref238}.

切片 Wasserstein 距离由文献\cite{ref266}引入，是为理解“迭代分布传递”（iterative distribution transfer）算法\cite{ref258}而提出的工具；其历史见 Bonnotte 的博士论文\cite{ref49}. 更多讨论见\cite[第5.5.4节]{ref281}，具有相似度量性质和统计性质的推广见\cite{ref233}. \autoref{prop:2.26}的推广见\cite{ref222}. 切片 Wasserstein 距离快速估计速率成立的尖锐条件由文献\cite{ref41}得到. 同一方向的其他结果，包括分布极限，见\cite{ref151,ref222,ref241,ref245,ref254,ref325,ref326}.
